\documentclass[11pt,a4paper]{article}
\usepackage{auditstyle}
\lhead{Fresh Glimm--Jaffe proof audit}
\rhead{16 August 2026}
\usepackage{bm}
\newcommand{\eps}{\varepsilon}
\newcommand{\T}{\mathbb T}
\newcommand{\dd}{\,d}
\DeclareMathOperator{\Tr}{Tr}

\title{Part II V1 in Lamport notation\\[3pt]
\large Thermal Weyl functionals, Mehler bridges, and the exact mean shift}
\author{Fresh reconstruction of \texttt{part\_ii/v1\_thermal\_representations.tex}}
\date{16 August 2026}

\begin{document}
\maketitle

\begin{resultbox}
All finite-dimensional identities in V1 are reconstructed below.  The reference
measure in the mean-shift formula is the \emph{free periodic Gaussian} measure.  The
continuum interacting marked-slice identity is not assumed here; it is proved by the
V3 reconstruction and consumed only afterward.
\end{resultbox}
\tableofcontents

\section{Objects and statement}

Fix $L,m>0$, $\eps_N=L/r_N$, the spatial lattice $\Lambda_N$ with $2r_N$ sites,
and
\[
 \Gamma_N=\{\pi j/L:-r_N\le j<r_N\},\qquad
 \gamma_N(k)^2=m^2+4\eps_N^{-2}\sin^2(\eps_Nk/2).
\]
On $L^2(\R^{\Lambda_N},d\bm\varphi)$ let
$\Phi_N(x)=\varphi_x$ and $\Pi_N(x)=-i\eps_N^{-1}\partial_{\varphi_x}$.  For real
$q,p$ set
\[
 \Phi_N(q)=\eps_N\sum_xq_x\Phi_N(x),\quad
 \Pi_N(p)=\eps_N\sum_xp_x\Pi_N(x),\quad
 W_N(\eps_Nq+ip)=e^{i(\Phi_N(q)+\Pi_N(p))}.
\]
Let $H_N^{(0)}$ be the nearest-neighbour oscillator Hamiltonian and
\[
 K_N=H_N^{(0)}+V_N,\qquad
 V_N=\lambda\eps_N\sum_x{:}\Phi_N(x)^4{:}_{c_{N,L}}.
\]
The model-sheet module proves that $K_N$ is self-adjoint, bounded below, has compact
resolvent, and has a simple ground state.

\section{Structured proof}

\LS{1}{1}{Derive the exact Schr\"odinger action of the Weyl operator.}

\LS{2}{1}{For $s\in\R$ define}
\[
 (U_s\psi)(\bm\varphi)=
 e^{is\eps_Nq\cdot\bm\varphi+\frac i2s^2\eps_Nq\cdot p}
 \psi(\bm\varphi+sp).
\]

\LS{3}{1}{$U_s$ is unitary.}

\Because Translation preserves Lebesgue measure and the prefactor has modulus one.

\LS{3}{2}{$U_{s+s'}=U_sU_{s'}$.}

\Because Direct substitution gives the exponent
\[
 s\eps_Nq\cdot\bm\varphi+s'\eps_Nq\cdot(\bm\varphi+sp)
 +\tfrac12(s^2+s'^2)\eps_Nq\cdot p
 =(s+s')\eps_Nq\cdot\bm\varphi+\tfrac12(s+s')^2\eps_Nq\cdot p.
\]

\LS{3}{3}{$s\mapsto U_s$ is a strongly continuous unitary group and its generator
extends $\Phi_N(q)+\Pi_N(p)$ on $C_c^\infty$.}

\Because Strong continuity first holds on $C_c^\infty$ by dominated convergence and
then on all $L^2$ by density and unitarity.  On the core,
\[
 \left.\frac d{ds}U_s\psi\right|_{s=0}
 =(i\eps_Nq\cdot\bm\varphi+p\cdot\nabla)\psi
 =i(\Phi_N(q)+\Pi_N(p))\psi.
\]
Stone's theorem gives a self-adjoint generator extending the displayed differential
operator.  Moreover $U_sC_c^\infty\subset C_c^\infty$ for every $s$; Nelson's
invariant-domain criterion therefore makes $C_c^\infty$ a core for that generator.
Thus the differential operator is essentially self-adjoint there and its closure
generates $U_s$.

\LS{2}{2}{At $s=1$,}
\begin{align}
 (W_N(\eps_Nq+ip)\psi)(\bm\varphi)
 &=e^{i\eps_Nq\cdot\bm\varphi+\frac i2\eps_Nq\cdot p}
   \psi(\bm\varphi+p),\tag{1.1}\\
 W_N(\eps_Nq+ip)
 &=e^{\frac i2\eps_Nq\cdot p}e^{i\Phi_N(q)}e^{i\Pi_N(p)}.\tag{1.2}
\end{align}

\LS{2}{3}{This proves $\langle1\rangle1$.}\QedStep

\LS{1}{2}{Prove trace class and the marked-slice kernel identity.}

\LS{2}{1}{$K_N$ has compact resolvent and is bounded below; hence
$e^{-tK_N}$ is trace class for $t>0$.}

\Because The model sheet proves
$K_N\ge a(-\Delta_{\bm\varphi}+|\bm\varphi|^2)-b$ in quadratic-form sense
for some $a>0,b<\infty$.  The comparison oscillator has compact resolvent and
summable heat eigenvalues; the min--max principle yields
$\Tr e^{-tK_N}<\infty$.

\LS{2}{2}{Let $k_t(\bm x,\bm y)$ be the continuous kernel of $e^{-tK_N}$.  For
bounded multiplication $f$ and translation $T_p\psi(x)=\psi(x+p)$,}
\[
 \Tr(e^{-tK_N}M_fT_p)=\int_{\R^{\Lambda_N}}
 k_t(\bm x+p,\bm x)f(\bm x)\,d\bm x.\tag{2.1}
\]

\Because Expand the trace in a real eigenbasis $K_Nu_j=E_ju_j$.  Its scalar series
is absolutely convergent because $M_fT_p$ is bounded.  To justify exchanging the
series and spatial integration, put $D(x)=e^{-(t/2)K_N}(x,x)$ and use
Cauchy--Schwarz twice:
\[
 \sum_j e^{-tE_j}|u_j(x)u_j(x+p)|
 \le \left(\sum_je^{-tE_j/2}\right)D(x)^{1/2}D(x+p)^{1/2},
\]
\[
 \int D(x)^{1/2}D(x+p)^{1/2}dx
 \le\left(\int D(x)dx\right)^{1/2}
     \left(\int D(x+p)dx\right)^{1/2}
 =\Tr e^{-(t/2)K_N}<\infty.
\]
Thus Fubini applies.  The kernel series is the continuous heat kernel, and the
change of variable in the trace gives (2.1) with $f(x)$ (not $f(x+p)$).

\LS{2}{3}{With $f(x)=e^{i\eps_Nq\cdot x}$, (1.1) and (2.1) give}
\[
 \Tr(e^{-tK_N}W_N(\eps_Nq+ip))
 =e^{\frac i2\eps_Nq\cdot p}\int
 e^{i\eps_Nq\cdot x}k_t(x+p,x)\,dx.\tag{2.2}
\]

\LS{2}{4}{This proves $\langle1\rangle2$.}\QedStep

\LS{1}{3}{Diagonalize the free system and compute its thermal Weyl functional.}

\LS{2}{1}{Set $X_x=\eps_N^{1/2}\Phi_N(x)$ and
$P_x=\eps_N^{1/2}\Pi_N(x)$.  The unitary discrete Fourier transform gives}
\[
 H_N^{(0)}=\sum_{k\in\Gamma_N}\gamma_N(k)(a_k^*a_k+\tfrac12),
\]
\[
 \widetilde X(k)=(2\gamma_N(k))^{-1/2}(a_k+a_{-k}^*),\quad
 \widetilde P(k)=-i(\gamma_N(k)/2)^{1/2}(a_k-a_{-k}^*).
\tag{3.1}
\]

\Because Fourier transformation diagonalizes the circulant nearest-neighbour matrix;
substitution of (3.1) and $[a_k,a_{k'}^*]=\delta_{kk'}$ cancels the $aa$ and
$a^*a^*$ terms and leaves the displayed oscillator sum.

\LS{2}{2}{For real $q,p$,}
\[
 \Phi_N(q)+\Pi_N(p)=\sum_k(\bar z_ka_k+z_ka_k^*),
\quad
 z_k=\eps_N^{1/2}\left[\frac{\widetilde q(k)}{\sqrt{2\gamma_N(k)}}
 +i\sqrt{\frac{\gamma_N(k)}2}\widetilde p(k)\right].\tag{3.2}
\]
Moreover
\[
 \sum_k|z_k|^2=\frac1{2(2r_N)}\sum_k
 \left(\gamma_N(k)^{-1}|\widehat q(k)|^2+\gamma_N(k)|\widehat p(k)|^2\right).
\tag{3.3}
\]

\Because Parseval gives (3.2).  In $|z_k|^2$ the cross term is
$\operatorname{Im}(\widetilde q(k)\overline{\widetilde p(k)})$; it is odd under
$k\mapsto-k$ and cancels in the sum.  The relation between the two Fourier
normalizations gives (3.3).

\LS{2}{3}{For one oscillator $H=\gamma(a^*a+1/2)$ and
$D(z)=e^{za^*-\bar za}$,}
\[
 \frac{\Tr(e^{-tH}D(z))}{\Tr e^{-tH}}
 =\exp[-\tfrac12|z|^2\coth(t\gamma/2)].\tag{3.4}
\]

\Because BCH gives $D(z)=e^{-|z|^2/2}e^{za^*}e^{-\bar za}$.  In the number
basis,
\[
 \langle n|e^{za^*}e^{-\bar za}|n\rangle
 =L_n(|z|^2),\qquad
 \sum_{n\ge0}w^nL_n(x)=\frac{e^{-xw/(1-w)}}{1-w}.
\]
The first identity follows by expanding both exponentials and matching the number
change; the second follows from
$\sum_{n\ge j}\binom njw^n=w^j(1-w)^{-j-1}$.  Put $w=e^{-t\gamma}$ and use
$(1+w)/(1-w)=\coth(t\gamma/2)$.

\LS{2}{4}{The finite tensor-product trace therefore equals}
\[
 F_N^{(0)}(t;\eps_Nq+ip)
 =\exp[-\tfrac14Q_N^{(t)}(q,p)],\tag{3.5}
\]
\[
 Q_N^{(t)}(q,p)=\frac1{2r_N}\sum_{k\in\Gamma_N}
 \coth(t\gamma_N(k)/2)
 \left(\gamma_N(k)^{-1}|\widehat q(k)|^2+\gamma_N(k)|\widehat p(k)|^2\right).
\tag{3.6}
\]
Also
$\Tr e^{-tH_N^{(0)}}=\prod_k[2\sinh(t\gamma_N(k)/2)]^{-1}$.

\Because Apply (3.4) to each commuting oscillator and insert (3.3); the zero-point
geometric sum is $[2\sinh(t\gamma/2)]^{-1}$.

\LS{2}{5}{This proves $\langle1\rangle3$.}\QedStep

\LS{1}{4}{Verify the Mehler kernel and the unique harmonic jump interpolant.}

\LS{2}{1}{For $H_\omega=\tfrac12(-\partial_x^2+\omega^2x^2)$, set}
\[
 k_\Delta^{\omega}(x,y)=
 \left(\frac{\omega}{2\pi\sinh(\omega\Delta)}\right)^{1/2}
 \exp\left[-\frac{\omega\coth(\omega\Delta)}2(x^2+y^2)
 +\frac{\omega xy}{\sinh(\omega\Delta)}\right].\tag{4.1}
\]

\LS{3}{1}{$\partial_\Delta k_\Delta^\omega=-H_\omega k_\Delta^\omega$.}

\Because Differentiate the prefactor and exponent, using
$(\omega\coth\omega\Delta)'=-\omega^2\operatorname{csch}^2(\omega\Delta)$ and
$(\omega\operatorname{csch}\omega\Delta)'=-\omega^2\coth(\omega\Delta)
\operatorname{csch}(\omega\Delta)$; the coefficients of
$1,x^2,xy,y^2$ agree separately.

\LS{3}{2}{$k_\Delta^\omega\to\delta_{x=y}$ as $\Delta\downarrow0$.}

\Because In variables $u=(x+y)/\sqrt2$, $v=(x-y)/\sqrt2$, the quadratic form
has coefficients $\omega\tanh(\omega\Delta/2)$ and
$\omega\coth(\omega\Delta/2)$; the latter tends to infinity with the normalized
Gaussian mass, while the former tends to zero.

\LS{3}{3}{Thus (4.1) is the heat kernel.}

\Because For compactly supported smooth initial data, both convolution with (4.1)
and $e^{-\Delta H_\omega}$ solve the same $L^2$ Cauchy problem.  Their difference
$v$ obeys $d\|v\|_2^2/d\Delta=-2\langle v,H_\omega v\rangle\le0$ and starts at zero.

\LS{2}{2}{For a prescribed jump $\rho$, the boundary-value problem}
\[
 -\ddot H+\omega^2H=0,\quad H(t)-H(0)=\rho,\quad \dot H(t)=\dot H(0)
\]
has the unique solution
\[
 H(s)=\frac\rho2\frac{\sinh(\omega(s-t/2))}{\sinh(\omega t/2)}.
\tag{4.2}
\]

\Because The general solution is $A\cosh(\omega s)+B\sinh(\omega s)$.  Solving the
two linear boundary equations gives $A=-\rho/2$ and
$B=(\rho/2)\coth(\omega t/2)$, which is (4.2) by the subtraction formula.
The homogeneous system has only $A=B=0$.

\LS{2}{3}{The interpolant satisfies}
\[
 |H(s)|\le|\rho|/2,\quad
 S(H):=\frac12\int_0^t(\dot H^2+\omega^2H^2)ds
 =\frac{\omega\rho^2}{4}\coth(\omega t/2).\tag{4.3}
\]

\Because $|s-t/2|\le t/2$ and $|\sinh|$ is increasing on $[0,\infty)$.
Integration by parts and the ODE give
$S(H)=\frac12[H\dot H]_0^t=\frac12\rho\dot H(0)$; differentiating (4.2) at
$0$ gives the last expression.

\LS{2}{4}{This proves $\langle1\rangle4$.}\QedStep

\LS{1}{5}{Prove the exact discrete Gaussian shift identity.}

Fix the equal partition $s_i=i\Delta$, $\Delta=t/M$, and let
$H_i=H(s_i)$.  Define the periodic Mehler-chain integral
\[
 J_0[f_0,\ldots,f_{M-1}]
 =\int\prod_{i=0}^{M-1}k_{\Delta}^\omega(\eta_{i+1},\eta_i)
 f_i(\eta_i)\,\prod_i d\eta_i,\qquad \eta_M=\eta_0,
\]
and the jump chain $J_\rho$ by replacing the closing condition by
$x_M=x_0+\rho$.

\LS{2}{1}{Under $x_i=\eta_i+H_i$, every interior Mehler exponent has zero linear
term in $\eta_i$.}

\Because Apart from its normalization, the negative logarithm of one kernel is
\[
 Q_\Delta(y,x)=\frac{\omega}{2\sinh(\omega\Delta)}
 [\cosh(\omega\Delta)(x^2+y^2)-2xy].
\]
For $1\le i<M$, the coefficient of $\eta_i$ in
$\sum_iQ_\Delta(\eta_{i+1}+H_{i+1},\eta_i+H_i)$ is proportional to
\[
 2\cosh(\omega\Delta)H_i-H_{i-1}-H_{i+1}=0,
\]
the exact three-term propagation law of every solution of
$-\ddot H+\omega^2H=0$.  The combined coefficient of the periodic variable
$\eta_0=\eta_M$ is the difference of the endpoint classical momenta and vanishes
because $\dot H(t)=\dot H(0)$.  Hence all linear terms, including the closing
one, vanish.

\LS{2}{2}{The constant part of the exponent is $S(H)$, independent of the mesh.}

\Because $Q_\Delta(H_{i+1},H_i)$ is exactly the classical action
$\frac12\int_{s_i}^{s_{i+1}}(\dot H^2+\omega^2H^2)ds$ of the unique harmonic
path joining the two endpoints.  The restriction of (4.2) is that path.  Hence
\[
 \sum_iQ_\Delta(H_{i+1},H_i)
 =\frac12\int_0^t(\dot H^2+\omega^2H^2)ds
 =\frac12[H\dot H]_0^t=S(H).
\]

\LS{2}{3}{The endpoint conditions transform correctly:}
$x_M=x_0+\rho$ becomes $\eta_M=\eta_0$ because $H(t)-H(0)=\rho$.

\LS{2}{4}{Therefore}
\[
 J_\rho[f_0,\ldots,f_{M-1}]
 =e^{-S(H)}J_0[f_0(\cdot+H_0),\ldots,f_{M-1}(\cdot+H_{M-1})].
\tag{5.1}
\]
The multimode identity is the product of (5.1).

\LS{2}{5}{This proves $\langle1\rangle5$.}\QedStep

\LS{1}{6}{Derive the interacting marked-slice mean-shift formula.}

\LS{2}{1}{The harmonic-reference FKN formula gives}
\[
 k_t^K(x,y)=k_t^{(0)}(x,y)
 \mathbb E_{x,y}^{t,0}\exp\left[-\int_0^tV_N(B_s)ds\right].\tag{6.1}
\]

\Because For bounded truncations of $V_N$ this is the Trotter product formula with
the Mehler kernel.  The Wick quartic is continuous and bounded below; truncate it
from above, apply monotone/dominated convergence to the positive kernels, and use
the closed lower-bounded form limit.  This proves (6.1) without importing an
interacting scaling limit.

\LS{2}{2}{Insert (6.1) into (2.2) and approximate the time integral by Riemann
sums.}

\Because Each bridge path is continuous and $V_N$ is continuous, so the Riemann sums
converge pointwise.  If $V_N\ge-c_N$, every exponential is bounded by $e^{tc_N}$;
the remaining Mehler chain has finite mass.  Dominated convergence applies.

\LS{2}{3}{Apply (5.1) mode by mode with jump
$\rho_k=\eps_N^{1/2}\widetilde p(k)$.}

\Because The position translation in (1.1) is $p$, hence in the rescaled mode
coordinate it is exactly $\rho_k$.  The product Jacobian is one.

\LS{2}{4}{The phase from (1.1) cancels the value of the marked position source on
$H(0)=-p/2$.}

\Because The inserted factor becomes
\[
 e^{\frac i2\eps_Nq\cdot p}e^{i\eps_Nq\cdot(\eta(0)-p/2)}
 =e^{i\eps_Nq\cdot\eta(0)}.
\]

\LS{2}{5}{Let $\mu_{N,t}^{\rm per}$ be the free periodic Gaussian field with mode
covariance}
\[
 c_\omega(s)=\frac{\cosh(\omega(t/2-|s|_{\T_t}))}
 {2\omega\sinh(\omega t/2)}.
\tag{6.2}
\]
Then
\[
 \frac{\Tr(e^{-tK_N}W_N(\eps_Nq+ip))}{\Tr e^{-tK_N}}
 =e^{-S_N(H_p)}
 \frac{\mathbb E_{\mu_{N,t}^{\rm per}}
 [e^{i\eps_Nq\cdot\eta(0)}e^{-\mathcal U_N(\eta+H_p)}]}
 {\mathbb E_{\mu_{N,t}^{\rm per}}[e^{-\mathcal U_N(\eta)}]}.
\tag{6.3}
\]

\Because The periodic Mehler chain is centered Gaussian and its two-point function is
(6.2), obtained either by inverting $-\partial_s^2+\omega^2$ on the time circle or by
the oscillator trace.  Normalize the periodic chain by the zero-source partition
function.  Steps $\langle2\rangle2$--$\langle2\rangle4$ give exactly the numerator;
$q=p=0$ gives the denominator.

\LS{2}{6}{The shifted Wick quartic is exactly}
\[
 {:}(\eta+H)^4{:}_c-{:}\eta^4{:}_c
 =4H{:}\eta^3{:}_c+6H^2{:}\eta^2{:}_c+4H^3\eta+H^4.
\tag{6.4}
\]

\Because Expand $(\eta+H)^4-6c(\eta+H)^2+3c^2$ and use
$\eta^3={:}\eta^3{:}_c+3c\eta$ and $\eta^2={:}\eta^2{:}_c+c$; all additional
$c$ terms cancel.

\LS{2}{7}{This proves $\langle1\rangle6$.}\QedStep

\LS{1}{7}{Construct the continuum free thermal functional.}

Let $\gamma=(-\partial_x^2+m^2)^{1/2}$ on $\T_L$ and
\[
 h_{\rm ct}=H^{-1/2}(\T_L)\oplus iH^{1/2}(\T_L).
\]

\LS{2}{1}{For $\zeta=q+ip\in h_{\rm ct}$, the symplectic form
$\operatorname{Im}\int\bar\zeta\zeta'$ is finite and defines the continuum Weyl
operator $W_{\rm ct}(\zeta)$.}

\Because Sobolev duality pairs $H^{-1/2}$ with $H^{1/2}$.  The covariance norm
$\langle q,\gamma^{-1}q\rangle+\langle p,\gamma p\rangle$ is finite by definition.

\LS{2}{2}{$e^{-t\,d\Gamma(\gamma)}$ is trace class and}
\[
 \Tr e^{-t\,d\Gamma(\gamma)}=\prod_{k\in(\pi/L)\mathbb Z}(1-e^{-t\gamma(k)})^{-1}.
\]

\Because The occupation basis gives the product of geometric series.  It converges
because $\sum_ke^{-t\gamma(k)}<\infty$.

\LS{2}{3}{For square-summable displacement amplitudes, absolute convergence permits
factorization of the trace into one-mode traces.  Hence}
\[
 F_\infty^{(0)}(t;q+ip)
 =\exp\left[-\frac14\sum_k\coth(t\gamma(k)/2)
 \left(\gamma(k)^{-1}|\widehat q(k)|^2+\gamma(k)|\widehat p(k)|^2\right)\right],
\tag{7.1}
\]
with the Fourier normalization fixed by the model sheet.

\Because $|\langle n,D(z)n\rangle|\le1$ dominates the trace by the partition
function; Fubini applies, and (3.4) evaluates every factor.

\LS{2}{4}{The D4 image of every fixed coarse symbol belongs to $h_{\rm ct}$.}

\Because Its Fourier coefficients contain
$P_\infty(\eps_Mk)$ with
$|P_\infty(\ell)|^2\le C_D(1+|\ell|)^{-2-\eta}$, $\eta>0$.  The momentum
sum is bounded by $\sum_k|k|(1+\eps_M|k|)^{-2-\eta}<\infty$, and the position
sum is easier.

\LS{2}{5}{This proves $\langle1\rangle7$.}\QedStep

\LS{1}{8}{Prove the free thermal fixed-coarse rate uniformly for $t\ge t_0>0$.}

\LS{2}{1}{For $w_t(\gamma)=\coth(t\gamma/2)$ and
$\bar w=\coth(t_0m/2)$,}
\begin{align*}
 1&\le w_t(\gamma)\le\bar w,\\
 |(\gamma w_t(\gamma))'|&\le\bar w+s_*,\\
 |(\gamma^{-1}w_t(\gamma))'|&\le(\bar w+s_*)m^{-2},
\end{align*}
where $s_*=\sup_{x\ge t_0m/2}x/\sinh^2x<\infty$.

\Because Differentiate $\coth(t\gamma/2)$ and use $\gamma\ge m$.

\LS{2}{2}{On low momenta, split each weighted multiplier difference into a
dispersion term and a D4-filter term.}

\Because For the momentum multiplier,
\[
 |\gamma_Nw_t(\gamma_N)|P_r|^2-\gamma w_t(\gamma)|P_\infty|^2|
 \le |\gamma_Nw_t(\gamma_N)-\gamma w_t(\gamma)|\,|P_r|^2
 +\gamma w_t(\gamma)||P_r|^2-|P_\infty|^2|.
\]
The first term is at most
$(\bar w+s_*)\eps_N^2k^4/(24m)$ times $|P_r|^2$; the position multiplier is
the same with the derivative bound $m^{-2}(\bar w+s_*)$.  The second terms use
the explicit finite-product filter tail.

\LS{2}{3}{On high momenta, multiply the vacuum D4 envelope estimate by at most
$\bar w$.}

\Because $1\le w_t\le\bar w$ uniformly in $t\ge t_0$; no derivative is used in
the tail.

\LS{2}{4}{Choose the same split exponent as the free-rate module,
$2-5\alpha=\alpha\eta=\theta=2\eta/(5+\eta)$.  Then}
\[
 |Q_{M+r}^{(t)}(\xi)-Q_\infty^{(t)}(\xi)|
 \le C_{M,\xi,m,L,t_0}2^{-\theta r},\qquad t\in[t_0,\infty].\tag{8.1}
\]

\Because The low sum is $O(4^{-r}(1+J_r)^5)$ and the high tail is
$O(J_r^{-\eta})$.  Put $J_r\asymp2^{\alpha r}$; the two exponents become equal.
Every sum and counting constant is the explicit one proved in the free-rate module.

\LS{2}{5}{The characteristic functionals obey the same rate up to a factor $1/4$.}

\Because For $x,y\ge0$, $|e^{-x/4}-e^{-y/4}|\le|x-y|/4$.

\LS{2}{6}{This proves $\langle1\rangle8$.}\QedStep

\section{V1 conclusion and dependency interface}

The proved outputs are the exact lattice Weyl phase, trace-kernel formula, free thermal
closed forms, Mehler/jump shift, interacting finite-lattice representation (6.3), and
uniform free fixed-coarse rate (8.1).  Formula (6.3) uses only a finite-dimensional
FKN argument.  The analogous continuum interacting formula requires the continuum
Hamiltonian/periodic FKN construction; the V3 module proves it before V4 uses it.

\end{document}
